%-------------------------------------------------------------------------------
%	SECTION TITLE
%-------------------------------------------------------------------------------
\cvsection{Projets et évènements}


%-------------------------------------------------------------------------------
%	CONTENT
%-------------------------------------------------------------------------------
\begin{cventries}
%---------------------------------------------------------
  \cventry
    {}
    {Recherche de motif dans un réseau : problème de plongements et d'isomorphismes de graphes. Résolution dans des cas remarquables}
    {Lycée Pierre-de-Fermat, Toulouse} % Location
    {Septembre 2024 - Juin 2025}
    {
        \begin{cvitems}
            \item {Notions de centre, centroïdes pour les graphes connexes acycliques.}
            \item {Réduction du problème en passant par l'algèbre dans le cas des graphes trivalents pour une résolution en temps polynomial.}
            \item Lien vers \href{https://github.com/Aur3m/tipe2-pres/}{la présentation et le rapport}, \href{https://github.com/Aur3m/tipe2}{le code en C++} 
        \end{cvitems}
    }
    
  \cventry
    {} % Organisation
    {Utilisation d'une hybridation entre la Recherche Tabou et la Colonie de fourmis pour résoudre le problème de la tournée de véhicules} % Project
    {Lycée Henri Poincaré, Nancy} % Location
    {Septembre 2022 - Juin 2023} % Date(s)
    {
      \begin{cvitems} % Description(s) of project
        \item {Étude d'une hybridation entre une colonie de fourmis et une recherche tabou pour le problème de tournées de véhicules}
      \end{cvitems}
    }
    
  \cventry
    {Classé 63ème national} % Organisation
    {Régionales puis Finale Concours Prologin} % Project
    {Epita, Paris} % Location
    {Mai 2023} % Date(s)
    {
      \begin{cvitems} % Description(s) of project
        \end{cvitems}
    }

  \cventry
    {En équipe de trois, nous nous sommes classés 82ème sur 1150 dans le monde} % Organisation
    {Reply Code Challenge - Teen 2022} % Project
    {En Ligne} % Location
    {Mars 2022} % Date(s)
    {
      \begin{cvitems} % Description(s) of project
        \end{cvitems}
    }

  \cventry
    {Club de mathématiques de l'Institut Mathématiques de Marseille (I2M) - Équipe classée 6ème sur 12 équipes finalistes} % Organisation
    {Régionales puis Finale du Tournoi Français des jeunes mathématicien(nes)} % Project
    {Bibliothèque du C.I.R.M., Marseille} % Location
    {Janvier 2022 - Mai 2022} % Date(s)
    {
      \begin{cvitems} % Description(s) of project
        \item Initiation au métier de chercheur en mathématiques et en informatique
        \item Rédaction en LateX
        \item Présentation des résultats devant un public avec contrainte de temps, \href{https://github.com/Aur3m/tfjm_souvenir/blob/main/Sujets_Mas__tfjm2_finale.pdf}{Lien vers la présentation}
      \end{cvitems}
    }

  \cventry
    {Blog de la section O.I.B. du Lycée Saint Charles} % Organisation
    {Un article de vulgarisation en anglais que j'ai écrit} % Project
    {Marseille, France} % Location
    {Mars 2022} % Date(s)
    {
      \begin{cvitems} % Description(s) of project
      \item Le but était d'introduire l'algorithmique à mes camarades de 1ère  \href{https://oibstcharles.blogspot.com/2022/03/algorithmics-my-personal-experience.html}{Lien vers l'article}
    \end{cvitems}
    }

  \cventry
    {Développeur et Administrateur Système} % Organisation
    {Communauté en ligne "La Confrérie"} % Project
    {Marseille, France} % Location
    {Mars 2018 - Décembre 2021} % Date(s)
    {
      \begin{cvitems} % Description(s) of project
      \item Développement d'un site web, de mods et de bots discord
      \item Administration Système : gestion d'une infrastructure réseau sous Kubernetes et Docker
    \end{cvitems}
    }

    \cventry
    {} % Organisation
    {Échanges en immersion dans des écoles en Angleterre} % Project
    {Teignmouth, Bournemouth et Eastbourne UK} % Location
    {juillet 2016 puis juin-juillet 2017} % Date(s)
    {
    }

    \cventry
    {Lectures scientifiques et articles de recherche} % Title
    {Lectures personnelles} % Organisation/Context
    {} % Location (laisser vide)
    {2025 - 2026} % Date(s)
    {
      \begin{cvitems}
        \item \textbf{Logique :} \textit{Proofs and Types} (J.-Y. Girard, Y. Lafont, P. Taylor)
        \item \textbf{Quelques articles qui m'ont marqué :}
        \begin{itemize}
            \item \textit{A new short proof of Kneser's conjecture} (Joshua E. Greene, 2002).
            \item \textit{On the size of Kakeya sets in finite fields} (Zeev Dvir, 2008).
            \item \textit{Induced subgraphs of hypercubes and a proof of the Sensitivity Conjecture} (Hao Huang, 2019).
        \end{itemize}
      \end{cvitems}
    }

\end{cventries}
